%%%%%%%%%%%%
%
% $Autor: Wings $
% $Datum: 2019-03-05 08:03:15Z $
% $Pfad: SafetyGuidelines.tex $
% $Version: 4250 $
% !TeX spellcheck = en_GB/de_DE
% !TeX encoding = utf8
% !TeX root = manual 
% !TeX TXS-program:bibliography = txs:///biber
%
%%%%%%%%%%%%

\chapter{Safety Guidelines}

\section{Safety Guidelines}

The following safety guidelines are essential to ensure reliable, secure, and fault-free operation of the License Plate Detection System.  
They cover electrical, thermal, mechanical, and software safety aspects, tailored to embedded AI hardware (Jetson Nano, camera, power supply).

\subsection{Electrical Safety}

\begin{itemize}
	\item Always use a certified **5 V / 4 A USB-C power supply** with proper overcurrent and short-circuit protection.  
	\item Avoid powering the unit when connecting or removing peripherals (camera, HDMI, USB).  
	\item Make sure all cables are insulated and rated for the current; check for worn or damaged wires before use.  
	\item Use grounded power outlets or surge protectors to prevent voltage spikes or earth loops.  
	\item Respect polarity: reverse connection may permanently damage the board or peripheral components.
\end{itemize}

\subsection{Thermal and Mechanical Safety}

\begin{itemize}
	\item Maintain adequate ventilation. Do not cover or obstruct the aluminum case’s fins or openings, as the board uses passive cooling.  
	\item Operate within safe ambient temperature (e.g. 0 °C to 40 °C). In hotter environments, consider adding an auxiliary fan.  
	\item Ensure no heavy objects rest on the system. Avoid placing the unit near liquids or in high humidity environments.  
	\item Secure the camera on its stand firmly to avoid fall damage or misalignment.  
	\item Use anti-static precautions (e.g. wrist strap) when handling the Jetson Nano board to prevent ESD damage.
\end{itemize}

\subsection{Operational Safety}

\begin{itemize}
	\item Boot and shut down properly via software commands (e.g. \texttt{sudo shutdown}) before removing power.  
	\item Monitor system logs and error messages to detect sensor faults, temperature warnings, or abnormal behavior.  
	\item Do not overload the USB ports with high-power devices that exceed the board’s power budget.  
	\item In case of smoke, smell, or abnormal heating, power off immediately and inspect hardware before reuse.
\end{itemize}

\subsection{Power \& Inference Safety}

According to studies on AI inference energy use, the power consumption of embedded devices performing computer vision tasks is nontrivial. For example, a work estimating power usage of heterogeneous devices shows that using YOLO models on devices like Jetson Nano yields measurable power-per-frame metrics 
Thus:

\begin{itemize}
	\item Regularly monitor power draw using a multimeter or inline ammeter to ensure it remains within safe limits (e.g. under the adapter’s rated capacity).  
	\item Avoid leaving the system running under heavy load continuously without periodic rest or thermal check.  
\end{itemize}

\subsection*{Summary}

Adhering to these safety rules protects both the hardware and the environment it operates in.  
Electrical and thermal control prevents damage, while mechanical and operational precautions ensure reliable use.  
Combining these with measurements and monitoring helps maintain long-term stability and prevent accidents.




